\documentclass[12pt]{article}
\usepackage[utf8]{inputenc}
\usepackage[T1]{fontenc}
\usepackage{hyperref}
\title{Tzeltal, Bachaj\'on Language Resources for K-12 Curriculum Development}
\author{Compiled by Claude Sonnet 5 (claude-sonnet-5) for Andr\'as Kornai}
\date{}
\begin{document}
\maketitle

\section{Introduction}
Tzeltal of Bachaj\'on (ISO 639-3 tzb) is a Mayan language of Chiapas, Mexico. The identification given with the task lists the family as ``Bookkeeping'', which looks like a data error. The searches found mostly mid-twentieth-century bilingual materials and one modern multidialect dictionary.

\section{Primary Educational Resources}
Catalogue records seen in search results (Colegio de M\'exico repository) describe a 1966 workbook for the first Bachaj\'on primer, a companion workbook for the second primer, and a 1973 Spanish and Bachaj\'on Tzeltal phrase booklet. Their record notes that the language used in that edition is no longer current. I did not confirm a direct link for each item, so they are not listed with URLs.

\section{Linguistic and Reference Materials}
\begin{itemize}
\item Gerdel and Slocum (1965), \emph{Vocabulario tzeltal de Bachaj\'on}, as indexed in Glottolog: \url{https://glottolog.org/resource/reference/id/77607}
\item Polian's Tseltal-Spanish multidialectal dictionary (8,109 entries, 20 dialects); I did not confirm that Bachaj\'on is included: \url{https://dictionaria.clld.org/contributions/tseltal}
\item Tzeltal language, Wikipedia: \url{https://wikipedia.com/wiki/ISO_639:tzb}
\end{itemize}

\section{Religious and Community Resources}
None were found in the searches.

\section{Digital and Audio Resources}
\begin{itemize}
\item Endangered Languages Project resource listing: \url{https://endangeredlanguages.com/zh-hans/node/40390/resources}
\end{itemize}

\section{Curriculum Development Context}
The sources suggest that Tzeltal is used almost exclusively orally and that schools rarely use Tzeltal materials. This is a general statement about Tzeltal, not specific to Bachaj\'on. The available Bachaj\'on-specific materials are old, so a curriculum would need new or revised materials.
\end{document}